% -------------------------------------------------------------------------------------
% Einlesen der .sty-Dateien
% -------------------------------------------------------------------------------------
%  se-pa1-input-styles.tex
%
%  Joerg Baumgart 01.08.2011
%
%  Zusammenfassung und Konfiguration wichtiger Styles f\"ur die 
%  Erzeugung von Seminar-, Projekt- und Bachelorarbeiten
%
%
\documentclass[12pt,BCOR=10mm,headinclude=on,footinclude=off,bibliography=totoc]{scrreprt}
\usepackage[T1]{fontenc}
\usepackage[utf8]{inputenc}
\usepackage[ngerman]{babel} % Deutsche Einstellungen
\usepackage{lmodern}

\usepackage{tikz} % Graphikpaket, das zu pdfLaTeX kompatibel ist
\usepackage{xkeyval} % Definition von Kommandos mit mehreren optionalen Argumenten
\usepackage{listings} % Formatierung von Programmlistings
\usepackage{graphicx} % Einbinden von Graphiken
\usepackage{ifthen}
\usepackage{color}
\usepackage{slashbox} % Diagonalen in Tabellenfeldern
\usepackage{framed} % Erzeugung schwarzer Linien am linken Rand zur Hervorhebung von Textteilen
\usepackage{caption} % Korrektes Setzen einer mehrzeiligen float-Unterschrift bei neu definierten float-Umgebungen
%\usepackage{floatrow}

% Es wird jeweils die sty-Datei importiert und entsprechende Konfigurationseinstellungen werden vorgenommen
%
\usepackage{sty/se-jb-scrpage2} % Formatierung der Kopf- und Fu{\ss}zeilen
\usepackage{sty/se-jb-footmisc}    % Fussnoten besser formatieren

\usepackage{sty/se-jb-glossaries} % Abk\"urzungsverzeichnis, Symbolverzeichnis, Glossar
   
\usepackage{sty/se-jb-floatrow}    % Definition und Konfiguration von float-Umgebungen (figure, table, die neue programm-Umgebung)
% Achtung: se-jb-varioref muss nach se-jb-floatrow importiert werden; 
% andernfalls ist der counter programm f\"ur die labelformat-Anweisung noch nicht definiert   
\usepackage{sty/se-jb-varioref}   % Definition von Querverweisen
\usepackage{sty/se-jb-chngcntr}   % Kapitelweise oder globale Nummerierung von Abbildungen etc.
   
\usepackage{sty/se-jb-listen} % Definition neuer, besser formatierter Listen
\usepackage{sty/se-jb-wa-kommandos} % neue Kommandos f\"ur Seminar-, Projekt- und Bachelorarbeiten


% -------------------------------------------------------------------------------------
% Individuelle Konfiguration des Dokumentes
% -------------------------------------------------------------------------------------
%  Individuelle Konfiguration einer Projektarbeit
%
%
%
%

%
% Literaturverzeichnis
% 
%\usepackage{se-jb-jurabib-theisen} % Literaturverzeichnis gem\"ass den Vorgaben von Theisen aufbauen

\usepackage{bibgerm}
\usepackage{url}
\usepackage{booktabs}
\usepackage{array}

\providecommand{\seCite}[3]{\ifthenelse{\equal{#2}{}}{#1 \cite{#3}}{#1 \cite[#2]{#3}}}


\providecommand{\seFootcite}[3]{\ifthenelse{\equal{#2}{}}{\footnote{#1 \cite{#3}}}{\footnote{#1 \cite[#2]{#3}}}}


% Weitere Optionseinstellungen f\"ur das Koma-Script
%
% Zwischen Abs\"atzen einen Abstand von 0.5 \baselineskip erzeugen
\KOMAoption{parskip}{full}
%
% Vergleiche Duden "Gliederung von Nummern, S.111" 
% DIN 5008 anschauen, wenn sie neu ver\"offentlicht wurde
\KOMAoption{numbers}{noendperiod}
%
%



%  Voreinstellungen f\"ur floats
%  Durch die verwendeten Parameter wird die Wahrscheinlichkeit deutlich kleiner, 
%  dass Gleitobjekte (z. B. Abbildungen) ans Ende des Dokumentes verschoben 
%  werden; 
%  Achtung: clearpage erzwingt die Ausgabe von Gleitobjekten
%
\renewcommand{\topfraction}{1}  % Gleitobjekte d\"urfen eine Seite zu 100% belegen 
\renewcommand{\bottomfraction}{1} % Entsprechender Wert f\"ur den unteren Teil der Seite
\renewcommand{\textfraction}{0} % Eine Seite darf auch ohne Fliesstext existieren
%%%\renewcommand{\floatpagefraction}{1} % Bedeutung unklar, daher keine Ver\"anderung des Vorgabewertes 
                                                                        % von 0.5; eventuell bringt ein \"Anderung auf 1 etwas, wenn 
                                                                         % Probleme mit floats auftreten
                                                                         
                                                                         
                                                                         
% Konfiguration von Programm-Listings
% 
% Achtung: hier gibt es nahezu beliebig viele weitere Konfigurationm\"oglichkeiten; vgl. Paketdokumentation
%
\lstset{language=[R/3 6.10]ABAP,basicstyle=\ttfamily,keywordstyle=\color{blue},captionpos=b,aboveskip=0mm,belowskip=0mm,
          xleftmargin=0em}        
          
%
% Grundkonfiguration der Abs\"ande zwischen den Items der maximal f\"unf Verschachtelungsebenen der 
% neuen Listenumgebungen
%                                                                             
% Initialisierung der Abst\"ande zwischen den items f\"ur seList; Grundeinheit: 0.5\baselineskip; siehe se-jb-listen
\seSetlistbaselineskip{1}{0.75}{0.75}{0.75}{0.75}
% Initialisierung der Abst\"ande zwischen den items f\"ur seToplist; Grundeinheit: 0.5\baselineskip; siehe se-jb-listen
\seSettoplistbaselineskip{1}{0.75}{0.75}{0.75}{0.75}     


%
%  Konfiguration der verschiedenen Verzeichnisse
%
%  abstandEintrag: Wert wird mit \baselineskip multipliziert
%

%
%  Abbildungsverzeichnis
%
\seKonfigurationAbb[
%verzeichnisname=Abbildungsverzeichnis,
labeltextLinks=, % kein Text links;
%labeltextRechts=:,
labelbreite=1cm,
%labeleinzug=1cm,
%abstandEintrag=1,
newpage=ja,
%pnumwidth=20mm,
%dotsep=1000,
%tocrmarg=4.5cm,
%abstandVerzeichnis=-1mm
]

%
% LIstingverzeichnis
%
\seKonfigurationPrg[
%verzeichnisname=Listing-Verzeichnis,
labeltextLinks=,
%labeltextRechts=:,
labelbreite=1cm,
%labeleinzug=2cm,
%abstandEintrag=1,
newpage=ja,
%%pnumwidth=20mm,
%dotsep=1000,
%tocrmarg=4.5cm,
%abstandVerzeichnis=-10mm
]

%
% Tabellenverzeichnis
%
\seKonfigurationTab[
%verzeichnisname=Liste der Tabellen,
labeltextLinks=,
%labeltextRechts=:,
labelbreite=1cm,
%labeleinzug=0.5cm,
%abstandEintrag=1,
newpage=ja,
%pnumwidth=20mm,
%dotsep=1000,
%tocrmarg=4.5cm,
%abstandVerzeichnis=-10mm
]

%
% Abk\"urzungsverzeichnis
%
\seKonfigurationAbk[
%verzeichnisname=Liste der Abk\"urzungen,
%labelbreite=3cm,
%labeleinzug=0.5cm,
%abstandEintrag=1,
%newpage=ja,
%abstandVerzeichnis=-10mm
]

%
% Symbolverzeichnis
% 
\seKonfigurationSym[
%verzeichnisname=Liste der Symbole,
%labelbreite=4cm,
%labeleinzug=3.5cm,
%abstandEintrag=1,
newpage=ja,
%abstandVerzeichnis=-10mm
]

%
% Glossar
%
\seKonfigurationGlo[
%verzeichnisname=Glossar,
%abstandEintrag=0,
]



% (eventuelle) Neudefinition f\"ur die Unter-/\"Uberschriften von Abbildungen, Tabellen und Listings
%
%
%\renewcommand{\seCaptionNameAbbildung}{Abb.}
%\renewcommand{\seCaptionNameTabelle}{Tab.}
%\renewcommand{\seCaptionNameProgramm}{Prg.}


% % (eventuelle) Neudefinition f\"ur Querverweise innerhalb des Textes
%
%
%
%\renewcommand{\seQuerverweisSeite}{Seite}
%\renewcommand{\seQuerverweisAbbildung}{Abb.}
%\renewcommand{\seQuerverweisTabelle}{Tab.}
%\renewcommand{\seQuerverweisProgramm}{Prg.}
%\renewcommand{\seQuerverweisKapitel}{Kap.}
%\renewcommand{\seQuerverweisGleichung}{Gl.}

% Kommandos, die direkt nach \begin{document} ausgef\"uhrt werden m\"ussen
%
%
%
\AtBeginDocument{%
\renewcommand{\listfigurename}{\seAbbildungenVerzeichnisname}
\renewcommand{\listtablename}{\seTabellenVerzeichnisname}
\renewcommand{\figurename}{\seCaptionNameAbbildung}
\renewcommand{\tablename}{\seCaptionNameTabelle}
\pagenumbering{roman}
}
                                                              
                                                                         

% -------------------------------------------------------------------------------------
% Individuelle Definition von Abk\"urzungen, Symbolen und eventuell Glossareintr\"agen
% -------------------------------------------------------------------------------------
%--------------------------------------------------------------------------------------
% Trennungsregeln
%--------------------------------------------------------------------------------------
\hyphenation{prob-lem-los}
\hyphenation{Inter-banken-han-del}
\hyphenation{Business-Objects}
\hyphenation{Web-ser-vice}
\hyphenation{Web-ser-vi-ces}
\hyphenation{Dash-board}
\hyphenation{Dash-boards}
\hyphenation{Zu-gangs-be-schrän-kung-en}
\hyphenation{ge-star-tet}
\hyphenation{NSFR}
\hyphenation{LCR}
\hyphenation{Va-ri-a-blen-aus-prä-gung-en}
\hyphenation{Dia-gramm}
\hyphenation{An-wen-dung}
\hyphenation{An-wen-dung-en}
\hyphenation{An-for-de-rung}
\hyphenation{An-for-de-rung-en}
\hyphenation{Haupt-an-for-de-rung-en}
\hyphenation{tech-nisch-en}
\hyphenation{Funk-ti-o-na-li-tät}
\hyphenation{Or-ches-trie-rung}
\hyphenation{Xcel-si-us}
\hyphenation{Xcel-si-us-A-dap-ter}
\hyphenation{Ein-stel-lung}
\hyphenation{Ein-stel-lung-en}
\hyphenation{Li-qui-di-tät}
\hyphenation{Li-qui-di-täts-ri-si-ko}
\hyphenation{Li-qui-di-täts-ri-si-ko-ma-nage-ment}
\hyphenation{Dash-board}


%--------------------------------------------------------------------------------------
% Abkürzungen
%--------------------------------------------------------------------------------------
\newacronym{dbms}{DBMS}{Da\-ten\-bank\-ma\-nage\-ment\-sys\-tem}
\newacronym{sql}{SQL}{Structured Query Language}
\newacronym{MaRisk}{MaRisk}{Mindestanforderungen an das Risikomanagement}
\newacronym{LiqV}{LiqV}{Liquiditätsverordnung}
\newacronym{ngap}{NGAP}{Next Generation ABAP Plattform}
\newacronym{erp}{ERP}{Enterprise Resource Planning}
\newacronym{crm}{CRM}{Customer Relationship Management}
\newacronym{saplrm}{SAP\,\,LRM}{SAP Liquidity Risk Management}
\newacronym{http}{HTTP}{Hypertext Transfer Protocol}
\newacronym{xml}{XML}{Extensible Markup Language}
\newacronym{NetWeaver}{NetWeaver}{SAP NetWeaver Application Server}
\newacronym{xcelsius}{Xcelsius}{BusinessObjects Xcelsius}
\newacronym{biz}{BIZ}{Bank für Internationalen Zahlungsausgleich}
\newacronym{url}{URL}{Uniform Resource Locator}
\newacronym{wysiwyg}{WYSIWYG}{What You See Is What You Get}
\newacronym{sdk}{SDK}{Software Development Kit}
\newacronym{lcr}{LCR}{Liquidity Coverage Ratio}
\newacronym{nsfr}{NSFR}{Net Stable Funding Ratio}
\newacronym{rest}{REST}{Representational State Transfer}
\newacronym{ria}{RIA}{Rich Internet Application}
\newacronym{abap}{ABAP}{Advanced Business Application Programming}
\newacronym{bw}{BW}{Business Warehouse}
\newacronym{libor}{LIBOR}{London Inter Bank Offered Rate}

%--------------------------------------------------------------------------------------
% Glossareinträge
%--------------------------------------------------------------------------------------

\newglossaryentry{glos:netFramework}{
first=.NET Framework\textsuperscript{GL},
name=.NET Framework,
description={ Das .NET Framework ist eine Plattform von Microsoft, mit der Anwendungen für das Betriebssystem Microsoft Windows erstellt und ausgeführt werden können. Die wichtigsten Komponenten sind Klassenbibliotheken, zum Beispiel für die Entwicklung der Oberflächen, und die Common Language Runtime. Die Anwendungen können in verschiedenen Programmiersprachen geschrieben werden. Zu den unterstützten Sprachen zählen unter anderem C++ und C\#. Der Quellcode wird dann in die Common Intermediate Language compiliert. Diese Zwischensprache kann dann von der Common Language Runtime ausgeführt werden.\seFootcite{vgl.}{S.382ff}{Lou.10}}
}


\newglossaryentry{glos:interbankenhandel}{
first=Inter\-banken\-han\-del\textsuperscript{GL},
name=Interbankenhandel,
description={Der Interbankenhandel ist der Handel von Wertpapieren, Anlagen oder ähnlichem zwischen Banken. Synonym wird auch oft der Begriff Interbankenmarkt verwendet. Für die Zinssätze, mit denen Banken untereinander handeln, existieren anerkannte Referenzen, wie zum Beispiel die \gls{libor}. In Liquiditätsengpässen kann der Interbankenhandel eine wichtige Refinanzierungsrolle darstellen. Der Handel zwischen Banken hängt sehr stark von dem gegenseitigen Vertrauen ab.\seFootcite{vgl.}{S.145f}{Wil.10}}
}

\newglossaryentry{glos:bankrun}{
first=Bankenpanik\textsuperscript{GL},
name=Bankenpanik,
description={Eine Bankenpanik ist ein Ereignis, bei dem eine große Anzahl von Anlegern versucht, ihre Einlagen bei einer Bank abzuziehen. Der Grund kann zum einen in der Veröffentlichung von schlechten Ergebnissen der Bank und damit in einem Vertrauensverlust begründet sein, zum anderen aber auch rein spekulativ sein. Für die Bank besteht die Gefahr der Insolvenz. Im Englischen spricht man von einem Bank Run.\seFootcite{vgl.}{S.1f}{Sch.11}}
}

\newglossaryentry{glos:sqlscript}{
first=SQLScript\textsuperscript{GL},
name=SQLScript,
description={SQLScript ist eine Erweiterung der Abfragesprache SQL und wird in der Datenbank von SAP HANA verwendet. Mit Hilfe von SQLScript lässt sich Anwendungslogik in die Datenbank auslagern. Dazu wurden unter anderem Datentypen, Prozeduren und Kontrollstrukturen hinzugefügt.\seFootcite{vgl.}{S.9f}{SQLScript.11}}
}

\newglossaryentry{glos:ria}{
first=\gls{ria}\textsuperscript{GL},
name=\gls{ria},
sort=RIA,
description={Unter dem Begriff Rich Internet Application werden Webanwendungen bezeichnet, die in ihrer Funktionalität und ihrem Aussehen Desktopanwendungen ähneln. Erstmals eingeführt wurde der Begriff von Macromedia. Zwischen normalen Webanwendungen und RIA kann keine klare Grenze gezogen werden. Ein wichtiges Indiz für eine RIA ist der Einsatz von Technologien wie zum Beispiel Adobe Flash, Adobe Air oder Microsoft Silverlight.\footnote{\seCite{vgl.}{S.32f}{DD.08} \seCite{und}{S.3f}{Pfe.09}\newline
Adobe Flash - http://www.adobe.com/products/flashplayer.html \newline
Adobe Air - http://www.adobe.com/products/air.html \newline
Microsoft Silverlight - http://www.microsoft.com/silverlight/} }
}

\newglossaryentry{glos:bydesign}{
first=SAP Business ByDesign\textsuperscript{GL},
name=SAP Business ByDesign,
description={SAP Business ByDesign ist eine Anwendung von SAP für mittelständige Unternehmen. Zu dem Funktionsunfang gehört sowohl ein \gls{erp}- als auch eine \gls{crm}-Lösung. Die Besonderheit von SAP Business ByDesign ist, dass es auf Servern bei SAP betrieben wird und Kunden die Anwendung mieten und über das Internet konsumieren.}
}

\newglossaryentry{glos:lcr}{
first=\gls{lcr}\textsuperscript{GL},
name=\gls{lcr},
sort=LCR,
description={
Die \gls{lcr}, die auch als Mindestliquiditätsquote bezeichnet wird, ist eine requlatorische Vorgabe im Rahmen von Basel III. Durch die Einhaltung soll sichergestellt werden, dass eine Bank unerwartete Zahlungsmittelabflüsse über einen Zeitraum von 30 Tagen mit eigenen liquiden Aktiva ausgleichen kann.\seFootcite{vgl.}{S.4}{BIS.10}
}
}

\newglossaryentry{glos:nsfr}{
first=\gls{nsfr}\textsuperscript{GL},
name=\gls{nsfr},
sort=NSFR,
description={
Die \gls{nsfr} ist eine requlatorische Vorgabe im Rahmen von Basel III und soll eine mittel- und langfristige Refinanzierung von Banken fördern. Er besteht aus dem Verhältnis von verfügbaren und erforderlichen langfristiger Refinanzierung. Sie wird auch als strukturelle Liquiditätsquote beizeichnet.\seFootcite{vgl.}{S.28}{BIS.10}
}
}

\newglossaryentry{glos:sdk}{
first=\gls{sdk}\textsuperscript{GL},
name=\gls{sdk},
sort=SDK,
description={
Ein \gls{sdk} besteht aus einer Reihe von Anwendungen, mit deren Hilfe Anwendungen erstellt werden können. Meistens gehört zum Umfang eines \gls{sdk} auch eine Dokumentation der Programmierschnittstelle und eine Entwicklungsumgebung dazu.
}
}

\newglossaryentry{glos:atom}{
first=Atom Syndication Format\textsuperscript{GL},
name=Atom Syndication Format,
description={
Das Atom Syndication Format ist ein XML-basierendes Format für den plattformunabhängigen Austausch von Informationen. Es wird in der Regel für sich periodisch änderne Informationen, wie zum Beispiel Newsfeeds, genutzt.\seFootcite{vgl.}{}{ATOM.07}
}
}

\newglossaryentry{glos:tomcat}{
first=Apache Tomcat\textsuperscript{GL},
name=Apache Tomcat,
description={
Der Apache Tomcat ist ein Projekt der Apache Software Foundation, mit dem Java-Code auf Webservern ausgeführt werden kann. Dadurch können dynamische Webanwendungen entwickelt werden.
}
}

\newglossaryentry{glos:rest}{
first=\gls{rest}\textsuperscript{GL},
name=\gls{rest},
sort=REST,
description={
\gls{rest} ist ein Paradigma für die Umsetzung von Webservices. Dabei sieht es vor, dass eine Ressource eindeutig über einen Uniform Resource Identifier identifiziert werden kann. Diese Ressource kann in verschiedenen Repräsentationen abgerufen werden. Auf Ressourcen lassen sich festgelegte Operationen anwenden. Die Kommunikation findet immer zustandslos statt.\seFootcite{vgl.}{S.245f}{RRD.07}
}
}

\newglossaryentry{glos:actionscript}{
first=ActionScript\textsuperscript{GL},
name=ActionScript,
description={
ActionScript ist die Programmiersprache für die Entwicklung von Adobe Flex-basierenden Programmen. ActionScript unterstützt die objektorientierte Programmierung und basiert auf dem Sprachkern von JavaScript. Für die Entwicklung existiert mit dem Flex Builder eine integrierte Entwicklungsumgebung.\seFootcite{vgl.}{S.3f}{Moo.07}
}
}


\newglossaryentry{glo:abap}{
first=\gls{abap}\textsuperscript{GL},
name=\gls{abap},
sort=ABAP,
description={
\gls{abap} ist eine Programmiersprache der vierten Generation. Sie wird für einen Großteil der von SAP entwickelten Anwendungen verwendet. Eine Besonderheit von \gls{abap} ist die einfache Integration von \gls{sql}-Abfragen. Mit \gls{abap} Objects ist auch die objektorientierte Programmierung möglich.
}
} 

\begin{document}

\newcommand{\version}{0.1}

% -------------------------------------------------------------------------------------
% Erzeugung des Titelblatts
% -------------------------------------------------------------------------------------
\seTitelblattZweiteProjektarbeit[
firmenlogo=images/sap-logo,
firmenlogoDeltaX=-13,
firmenlogoDeltaY=-10,
thema=Entwicklung einer Zwischenschicht für die Nutzung weiterer Anwendungen in Verbindung mit der Berechnungskomponente des Liquidity Risk Managements,
verfasser=Fabian Kajzar,
matrikelnummer=428094,
kurs=WWI\,09\,SW\,B,
firma=SAP AG,
abteilung=Application Strategic Innovation - HPA,
studiengangsleiter=Prof. Dr.-Ing. J\"org Baumgart,
wissenschaftlicherBetreuerName=Prof. Dr. Hans-Henning Pagnia,
wissenschaftlicherBetreuerEmail=hans-henning.pagnia@dhbw-mannheim.de,
wissenschaftlicherBetreuerTelefon=0621 4105-1131,
firmenbetreuerName=Jens Mett,
firmenbetreuerEmail=jens.mett@sap.com,
firmenbetreuerTelefon=06227 7-61785,
bearbeitungszeitraumVon=13. Februar 2012,
bearbeitungszeitraumBis=4. Mai 2012,
sperrvermerk=nein
]

\pagenumbering{Roman}
\setcounter{page}{0}

% -------------------------------------------------------------------------------------
% Erzeugung der Kurzfassung; Verfasser, Firma und Thema werden automatisch \"ubernommen
% -------------------------------------------------------------------------------------
\seKurzfassung{}

\glsresetall

% Beispiel f\"ur ein Kapitel, dass vor dem Einleitungskapitel kommt, z. B. ein Vorwort oder eine Danksagung
%\seKapitelVorEinleitung{Vorwort}

% -------------------------------------------------------------------------------------
% Ausgabe des Inhaltsverzeichnisses
% -------------------------------------------------------------------------------------
\seInhaltsverzeichnis[
einrueckung=ja,
gliederungsebenen=4
]

% -------------------------------------------------------------------------------------
% Ausgabe der verschiedenen Verzeichnisse
% -------------------------------------------------------------------------------------
\seVerzeichnisse[gliederungsebene=section,imInhaltsverzeichnis=ja]{abk}{sym}{abb}{tab}{prg}


% -------------------------------------------------------------------------------------
% Vorbereitung für Inhalt der Arbeit
% -------------------------------------------------------------------------------------


% -------------------------------------------------------------------------------------
% Inhalt der Arbeit
% -------------------------------------------------------------------------------------
\chapter{Testkapitel}
\pagenumbering{arabic}

\section{Bild}

\begin{figure}[h]
\centering
\setlength{\unitlength}{1mm}
\includegraphics[width=15cm]{images/sap-logo.png}
\caption[Umfang der Cloud-Angebote als Venn-Diagramm]{Umfang der Cloud-Angebote als Venn-Diagramm, eigene Darstellung\label{fig:grundlagen:venn}}
\end{figure}

\section{Tabelle}

\begin{table}[htbp]%
\centering%
\begin{tabular}{| p{3cm} | p{5cm} |}
\hline
Speicher & Zugriffszeit \\
\hline
\hline
L1 Cache & ca. 4 CPU-Zyklen  \\ \hline
Hauptspeicher & 100 - 400 CPU-Zyklen \\ \hline
Festplatte & 1.000.000 CPU-Zyklen \\ \hline
\end{tabular} 
\caption[Vergleich der Zugriffszeiten von Speicher]{Vergleich der Zugriffszeiten von Speicher\label{table:zugriff}\protect \footnotemark}
\end{table}
\footnotetext{footnotetext...}

\section{Glossar}
Das hier ist ein Eintrag im Glossar: \gls{glos:test}

\section{Abkürzung und Symbol}
1. Verwendung einer Abkürzung \gls{test}
2. Verwendung der gleichen Abkürzung \gls{test}
Symbol: \gls{symb:pi}

\section{Listing}
Ein Beispielprogramm mit ABAP Syntax-Highlighting ist in \vref{hello} zu sehen.

\begin{programm}[htbp]
\begin{lstlisting}
REPORT TEST.

DATA LASTNAME TYPE STRING.

WRITE 'Hello World'.
REPLACE 'A' WITH 'B' INTO LASTNAME.

SELECT * FROM FLIGHTINFO WHERE CLASS = 'Y' OR CLASS = 'C'.
\end{lstlisting}
\caption{Die Klasse \texttt{HelloDHBW}\label{hello}}
\end{programm}

\section{Zitat}
Einfaches Zitat in Fußnote.\seFootcite{vgl.}{S.4}{Kaj.12a} Komplexer Zitierter Text in Fußnote.\footnote{\seCite{vgl.}{S.7}{Kaj.12b} \seCite{und}{S.4}{Gol.12}}
Noch ein Zitat\seFootcite{vgl.}{S.4}{KÜ.12}
Noch ein Zitat mit drei autoren\seFootcite{vgl.}{S.4}{KMS.08}

\section{Formel}
Text $\int\limits_a^b {\frac{d}{dx}F(x)dx} = F(b)-F(a)$ der inline mit einer Formel steht ..der inline mit einer Formel steht ..der inline mit einer Formel steht

Es gehen auch numerierte Formeln:
\begin{equation}
\int\limits_a^b {\frac{d}{dx}F(x)dx} = F(b)-F(a)
\end{equation}
die dann zentriert dargestellt weredn

\section{Footnote}
Kurzer Text vor der Fußnote.\footnote{eine ganz normale Fußnote, reiner Text...} Kurzer Text nach der Fußnote.


% -------------------------------------------------------------------------------------
% Anhang der Arbeit
% -------------------------------------------------------------------------------------
\seAppendix{}

\pagenumbering{Roman}
\setcounter{page}{7}

\chapter{Anhang}

Inhalt des Anhangs

% -------------------------------------------------------------------------------------
%  Erzeugung eines Glossars
% -------------------------------------------------------------------------------------
\newpage
\sePrintGlossary{}

% -------------------------------------------------------------------------------------
% Literaturverzeichnisses
% -------------------------------------------------------------------------------------
\sePrintBibliography{}

\bibliographystyle{mystyle}
\seBibliography{wa}

% -------------------------------------------------------------------------------------
% Erzeugung der ehrenw\"ortlichen Erkl\"arung
% -------------------------------------------------------------------------------------
\seEhrenwoertlicheErklaerung{}

\end{document}
